% Generated draft: requires mathematical and source review.
\documentclass[11pt]{article}
\usepackage[margin=25mm]{geometry}
\usepackage{fontspec}
\setmainfont{Latin Modern Roman}
\usepackage{amsmath,amssymb,mathtools,unicode-math}
\setmathfont{Latin Modern Math}
\usepackage{enumitem,needspace,xurl,hyperref}
\setlength{\parindent}{0pt}
\setlength{\parskip}{5pt}
\setlength{\emergencystretch}{3em}
\newcommand{\N}{\mathbb{N}}
\newcommand{\Z}{\mathbb{Z}}
\newcommand{\Q}{\mathbb{Q}}
\newcommand{\R}{\mathbb{R}}
\newcommand{\C}{\mathbb{C}}
\newcommand{\F}{\mathbb{F}}
\newcommand{\A}{\mathbb{A}}
\newcommand{\PP}{\mathbb{P}}
\newcommand{\E}{\mathbb{E}}
\newcommand{\eps}{\varepsilon}
\newcommand{\dd}{\,\mathrm{d}}
\begin{document}
\section*{UnsolvedMath ID: 20000304}
\textbf{Source problem code:} catalogue entry
\subsection{Definitions and mathematical statement}
Alexeev's question asks for definitions of log canonical, log terminal and related singularities for a pair \((X,Z)\), where \(X\) is normal and \(K_X\) need not be \(\Q\)-Cartier, and \(Z\) is a formal sum of subschemes. The original does not specify the field or allowed coefficients. The attached aid uses characteristic zero and an effective rational sum \(Z=\sum_jt_jZ_j\), with ideals \(\mathfrak a_j\ne0\) and \(t_j\in\Q_{\geq0}\).

For a proper birational morphism \(f:Y\to X\) with smooth \(Y\) and compatible canonical divisors, define the natural pullback of a Weil divisor \(D\) by
\[
\mathcal O_Y(-f^{\natural}D)=
(\mathcal O_X(-D)\cdot\mathcal O_Y)^{\vee\vee}.
\]
The asymptotic pullback \(f^*D\) is obtained coefficientwise from \(m^{-1}f^{\natural}(mD)\) along sufficiently divisible indices. For a prime divisor \(E\) over \(X\), put \(\operatorname{ord}_E Z=\sum_jt_j\operatorname{ord}_E\mathfrak a_j\). The limiting and plus discrepancy expressions are
\[
a^-_{m,E}(X,Z)=\operatorname{ord}_E\!\left(K_Y-\frac1m f^{\natural}(mK_X)\right)-\operatorname{ord}_E Z,
\]
\[
a^+_E(X,Z)=\operatorname{ord}_E\bigl(K_Y+f^*(-K_X)\bigr)-\operatorname{ord}_E Z.
\]
Their log-canonical and log-terminal tests use thresholds \(\geq-1\) and \(>-1\), respectively, for the relevant divisors, with the limiting definition also requiring a specified sufficiently-divisible-\(m\) convention. These competing extensions must be distinguished; in the \(\Q\)-Gorenstein case and divisible index, they recover the usual discrepancy expressions. The aid reports de Fernex--Hacon's 2009 limiting/plus theory as a published answer to the historical request, without supplying all definitional quantifiers for every ``etc.'' class.

Its rational compression result chooses \(q\geq1\) with \(qt_j\in\Z\) and sets
\[
\mathfrak b=\prod_j\mathfrak a_j^{qt_j},\qquad
\operatorname{ord}_E Z=\frac1q\operatorname{ord}_E\mathfrak b.
\]
It reports that the limiting/plus and Mather--Jacobian subscheme contributions depend only on this weighted ideal up to integral closure. This claim has the aid's rational effective-coefficient scope, not arbitrary real or signed sums.

The aid further reports the exact subscheme-independent gap
\[
a^+_E(X,Z)-a^-_{m,E}(X,Z)
=\operatorname{ord}_E\!\left(f^*(-K_X)+\frac1m f^{\natural}(mK_X)\right)\geq0.
\]
It says this gap decreases along divisible indices and, along factorial indices, converges to
\[
\operatorname{ord}_E\bigl(f^*(-K_X)+f^*K_X\bigr),
\]
the asymptotic sign-linearity defect. No equivalence between vanishing of this defect and every possible numerical or \(\Q\)-Cartier condition is asserted here.
\subsection{Short English statement}
Define singularity classes for normal pairs without a \(\Q\)-Cartier canonical divisor. The attached aid reports the existing limiting/plus discrepancy approach, compresses rational subscheme sums into one weighted product ideal, and compares the two discrepancies by a nonnegative sign-linearity defect.
\subsection{Sources}
\begin{itemize}
\item Local UnsolvedMath registry, permanent ID \(20000304\), exact public code \texttt{catalogue entry}; canonical location \texttt{aim-algebraic-geometry-notes.json}, notes[303].
\item AIM workshop ``Numerical invariants of singularities and higher-dimensional algebraic varieties,'' 31 July--4 August 2006, Problem 1, attributed to Alexeev; supplied URL \url{https://aimath.org/WWN/singularvariety/problems_final.pdf}. The supplied verification reports legibility/agreement with the PDF; this drafting step did not consult it. The original literature field is empty.
\item Attached research attempt 1, ``Rational ideal compression and the non-Q-Gorenstein sign defect,'' identifying de Fernex--Hacon's 2009 theory and mentioning Mather--Jacobian discrepancies without full publication citations.
\end{itemize}
\end{document}
